% ─────────────────────────────────────────────────────────────────────────────
\section{Summary}
\label{sec:summary}
% ─────────────────────────────────────────────────────────────────────────────

This study asked whether classical machine learning on flight-level statistical
features can diagnose aircraft engine health on the real-world NGAFID Cessna-172
benchmark, decomposed into binary Anomaly Detection (AD) and 19-class Fault
Classification (FC). The most important findings are:

\begin{itemize}
    \item \textbf{Anomaly Detection is well within reach of classical ML.}
          A tuned Ensemble on global statistics (Repr.~F\textsubscript{++})
          achieves F1\,=\,0.703 (AUC\,=\,0.767); on Repr.~B the same Ensemble
          reaches 0.691, ahead of Logistic Regression (0.679). The result is
          interpretable, laptop-trainable, and dominated by oil pressure---a
          physically sensible early health indicator.

    \item \textbf{Fault Classification is the hard, unsolved half.} Flat XGBoost
          reaches only F1\,=\,0.116 on Repr.~B (0.176 tuned); cross-cylinder
          features (Repr.~F\textsubscript{++}) raise the flat best to 0.207, and
          an Oracle K\,=\,4 hierarchical grouping lifts the \emph{ceiling} to
          0.459. Minority fault classes remain largely unrecognized under a 38:1
          imbalance.

    \item \textbf{The bottleneck is fault typing, not detection.} In the
          end-to-end cascade (F1\,=\,0.131), Stage~2 wrong-type errors are the
          single largest error category (20.0\% of predictions), so improving FC
          has more leverage than improving AD.

    \item \textbf{The remaining gap is temporal in nature.} Global and segmented
          statistics plateau because fault signatures are local micro-patterns
          living in low-variance principal components (PC9 and PC19--23,
          $<$0.1\% of variance). Locally restricted temporal deep learning is the
          indicated route to close it.

    \item \textbf{The deep learning comparison must be read with caution.} The
          benchmark figures used as an upper bound come from a preprint
          \cite{lmsd2025} that was later withdrawn by its authors over
          metric-computation flaws; the qualitative conclusions hold, but the
          exact gap magnitudes are provisional.
\end{itemize}

In short, classical ML delivers a deployable anomaly screen for general-aviation
condition-based maintenance today, while reliable fault attribution awaits
temporal models and a re-validated benchmark. All code is available at
\url{https://github.com/Barbosa12/NGAFID}.
